% chktex-file 18 disable warnings about "
% chktex-file 19 disable warnings about ô
% chktex-file 6 disable italic correction warnings


\chapter{Revisão sobre Redes Neurais}

As primeiras definições de uma rede neural artificial foram feitas em 1943 por
McCulloch e Pits inspiradas na hipótese de que as atividades mentais do cérebro
humano consistem em reações eletroquímicas em redes de células cerebrais
chamadas de neurônios. O neurônio artificial consiste em uma unidade
computacional abstrata capaz de receber estímulos de entrada e produzir uma ou
mais saídas e uma rede neural artificial --- de agora em diante referenciada
apenas por \emph{rede neural} --- consiste em um grafo direcionado cujos nodos
são neurônios artificiais.

\begin{figure}
  \caption{Abstração de um neurônio artificial}
  \begin{center}
    \includegraphics[scale=0.5]{placeholder}
  \end{center}
\end{figure}

\section{O neurônio artificial}

Cada neurônio recebe $m$ entradas ponderadas por um peso $w$,
correspondente a entrada do seu i-ésimo vizinho. Esta configuração foi
inspirada nos neurônios naturais cujas interconexões são reguladas por um
potencial de ativação responsável por diferenciar a intensidade entre as
mesmas. A saída é produzida aplicando-se a função limite $\sigma$ na soma
de todos as entradas ponderadas do neurônio. Esta função simula a \emph{lógica
  de limite} existente nos neurônios naturais, os quais somente irão propagar o
seu sinal caso haja um acúmulo suficiente de potencial elétrico no mesmo. Em
sua forma matemática, a saída $y$ do neurônio artificial pode ser expressa da
seguinte forma:

$$y=\sigma(\sum_{i=1}^{m}x_i w_i + b)$$

$\sigma$ representa uma classe de funções que podem ser utilizadas dependendo
do objetivo da construção da rede neural. Para problemas de classificação
\emph{linearmente separáveis}, pode ser usada a função limite exemplificada na
figura~\ref{fig:lim-fun}, enquanto para problemas não linearmente separáveis é
necessário utilizar algum tipo de função \emph{sigmoide}. A função logística é
um exemplo de função sigmoide frequentemente utilizada por ser diferenciável em
todo o seu domínio, o que permite produzir garantias de convergência nos
algoritmos de aprendizagem.

O valor $b$ corresponde ao \emph{bias} aplicado a função e pode ser omitido
simplesmente inserindo-se artificialmente uma nova entrada $x_0 = 1$.

\begin{figure}\label{fig:lim-fun}
  \caption{Exemplo de funções de limite $\sigma$ usadas em neurônios artificiais}
  \begin{center}
    \includegraphics[scale=0.5]{placeholder}
  \end{center}
\end{figure}

\subsection{Bias}

Todo neurônio artificial é dito ter pelo menos uma entrada extra $x_0 = 1$ e
consequentemente, o peso extra associado $w_0$ chamado de bias $b$. O bias
permite que a função aproximada pelo neurônio não necessite obrigatoriamente
passar pela origem $(0,0)$, o que reduz o erro médio da aproximação, caso os
dados de treinamento não tenham sofrido nenhum tipo de ajuste numérico, como por
exemplo, a remoção do valor médio.

\begin{figure}
  \caption{Exemplo de função aproximada com e sem o vetor de bias}
  \begin{center}
    \includegraphics[scale=0.5]{placeholder}
  \end{center}
\end{figure}

\subsection{Redes neurais}

Em uma rede neural, os neurônios são organizados em \emph{camadas}. Alguns
autores não consideram a camada de entrada como uma camada distinta, mas esta
abstração permite simplificar a notação e torna o processo de conectar redes
neurais entre si mais intuitivo. Outra camada especial é a de saída, pois é
ela que determina se a rede exercerá a função de um \emph{classificador} ou de
um \emph{regressor}. Todas as demais camadas intermediárias são chamadas de
\emph{camadas ocultas}; o nome não possui nenhuma conotação especial e serve
apenas para diferenciá-las das camadas de entrada e saída. Em uma rede
\emph{feed forward} as saídas de cada camada somente são computadas na direção
e sentido da camada de saída, isto é, não há uma retroalimentação das conexões
de saída para as de entrada. Quando isto acontece a rede é chamada de
\emph{recorrente} porém, o estudo deste tipo de rede neural não esta nos
objetivos deste trabalho.

Sendo assim, uma rede neural necessita ter pelo menos duas camadas --- a de
entrada e a de saída. A adição de camadas ou neurônios não possui um efeito
óbvio e generalizável para qualquer rede, e portanto é alvo de constante
experimentação e observação empírica sujeito a comportamento variável
dependendo do problema.

\begin{figure}\label{fig:003-nn}
  \caption{Exemplo de uma rede neural de 3 camadas}
  \begin{center}
    \includegraphics[scale=0.5]{placeholder}
  \end{center}
\end{figure}

Cada camada $S$ é formada por $R$ neurônios, cujos pesos das conexões podem ser
representados através de uma matriz $\mathbf{W}$. A notação em forma vetorial e
matricial dos pesos das camadas permite que operações mais eficientes sejam
utilizadas para realizar a computação da saída da rede neural e do ajuste dos
mesmos.

$$ W_{S,R} =
\begin{bmatrix}
  w_{1,1} & w_{1,2} & \cdots & w_{1,R} \\
  w_{2,1} & w_{2,2} & \cdots & w_{2,R} \\
  \vdots  & \vdots  & \ddots & \vdots  \\
  w_{S,1} & w_{S,2} & \cdots & w_{S,R}
\end{bmatrix}
$$

A saída ou \emph{ativação} de uma camada $\mathbf{\vec a^{\prime}}$ pode ser expressa em função da
ativação da camada anterior $\mathbf{\vec a}$ da seguinte maneira:

$$\mathbf{\vec a^{\prime}} = \sigma(\mathbf{\vec W\vec a+\vec b})$$

\subsection{Aprendizagem}

A aprendizagem em uma rede neural se da pelo ajuste dos pesos $w$ das conexões
dos neurônios. Ao longo dos anos foram propostas várias maneiras de atualizar
os pesos de uma rede neural automaticamente através de uma \emph{regra de
  aprendizagem}. Estas regras foram categorizadas em três grandes grupos:

\begin{description}

  \itemsep1pt\parskip0pt\parsep0pt

  \item[Aprendizagem Supervisionada] As entradas são acompanhadas por um
    conjunto de valores corretamente associados a mesma.

  \item[Aprendizagem Por Reforço] Semelhante ao aprendizado supervisionado,
    porém ao invés da resposta correta, é fornecido algo como uma nota ou uma
    indicação da performance da predição.

  \item[Aprendizagem Não Supervisionada] Nenhuma informação adicional a
    respeito das entradas é fornecida ao algoritmo.

\end{description}

\begin{figure}\label{fig:limit-functions-exapmple}
  \caption{Exemplo de funções de ativação}
  \begin{center}
    \includegraphics[scale=0.5]{placeholder}
  \end{center}
  \legend{Na esquerda uma simples função de limite; à direita a função $\frac{1}{1+e^{-x}}$}
\end{figure}

\subsubsection{Backpropagation}

Obviamente que a aplicabilidade de redes neurais seria extremamente limitada se
não houvesse uma maneira eficiente de aprender os pesos da rede de forma
automática.  Atualmente, a melhor maneira conhecida de fazer o mesmo é através
do algoritmo de \emph{backpropagation}. Este algoritmo embora já sabido desde a
década de 60, ganhou visibilidade na área de IA na década de 80, quando
aplicado em redes neurais.\citep{russell1995modern}

O processo de aprendizagem de uma rede neural segue a noção intuitiva que temos
do assunto: Um individuo faz uma predição acerca de uma proposição ou pergunta,
avalia o \emph{feedback} recebido do ambiente e corrige a sua previsão baseado
na distância entre o que foi previsto e o que de fato aconteceu.

A determinação da direção no espaço de pesos da correção é dado pelo
\emph{gradiente} função de erro $\mathbf{L}$, ou como é mais comumente
conhecida, \emph{loss function}. A eficiência do algoritmo de backpropagation
se baseia no fato de que a computação de cada camada depende somente da
computação das camadas anteriores.

Embora o algoritmo use a derivada da \emph{loss function} pré calculada de
maneira analítica, é comum ver nas implementações o uso da computação do
gradiente de maneira iterativa como forma de validação da implementação.

A computação completa da rede neural então, se da em dois passos: O primeiro
chamado de \emph{forward pass}, a ativação de cada camada é computada até que
se chegue numa resposta na camada de saída, a qual é usada para calcular o erro
em cada neurônio e o ajuste dos pesos $\Delta W_{i,j}$ do mesmo de acordo com a
seguinte equação:

$$ \mathbf{\Delta W_{i,j} = -\alpha\frac{\partial L}{\partial W_{i,j}}} $$

Onde $\alpha$ representa a taxa de aprendizagem.
